\section{Security Analysis}
\label{sec:security}

\todo{The statements below are proof sketches to be turned into full proofs
(or a mechanized model) before submission. Nothing here is claimed as proved yet.}

\subsection{Isolation outside misses}
\begin{proposition}[sketch]
\label{prop:oi}
Consider any two executions in which principal $p$ submits the same requests at
the same times and $\Lpub$ takes the same values. If no request of $p$ misses its
deadline in either execution, then $p$'s admission outcomes, first-token release
times and usage records are identical in both.
\end{proposition}

\noindent\emph{Sketch.} By induction over $p$'s requests. $p$'s \vcd list is
updated only by $p$'s own requests (\cref{sec:vcd}), and the public set is
fixed, so $\pv$ is the same in both executions. The deadline in
\cref{eq:deadline}, the admission test in \cref{alg:conductor} and the usage
record are functions of $(t_r, n, \pv, \Lpub)$ only. Without a miss, release
happens at $D_r$. Tokens after the first are released with their measured
spacing, which depends on decode-batch contention; we place that in $\Lpub$'s
scope (\cref{sec:threat}) and do not claim isolation for it. \qed

\subsection{Bounding what misses leak}
A miss releases the first token at $f_r > D_r$, and $f_r$ can depend on
physical state. Two factors limit the damage. First, misses are observable to
\sys, so they can be counted, rate-limited and reported. Second, after each
miss \sys enlarges the principal's bucket for the rest of the epoch, the
penalty schedule used in predictive mitigation, for which Askarov, Zhang and
Myers bound total leakage by a polylogarithmic function of elapsed
time~\cite{askarov2010predictive,zhang2011predictive}.

\begin{conjecture}
\label{conj:bound}
Over $E$ epochs with $M$ misses in total, a principal learns at most
$E\log_2 k$ bits from $\Lpub$ plus $O(M \log_2 (T_{\max}/\beta))$ bits from
misses, where $T_{\max}$ is the largest release delay.
\end{conjecture}
\todo{Prove by adapting the predictive-mitigation argument; check that its
assumptions (a public schedule, a bounded set of release times) hold for
\cref{eq:deadline}.}

\subsection{Coverage of the attacks}
\begin{itemize}
\item \textbf{\A1--\A3 (timing).} Global hits, the tier that served a block and
  its replica count change only $f_r$, never $D_r$. Outside misses they are
  invisible.
\item \textbf{\A4 (admission).} Admission depends on $\pv$ and $\Lpub$, not on the
  global match.
\item \textbf{\A5 (accounting).} Cached tokens are reported as $\pv$.
\item \textbf{\A6 (eviction).} The attacker's \emph{virtual} list is evicted only
  by its own requests. Physical eviction of its blocks causes recomputation and,
  at worst, a miss, which the pin budget makes rare and the bound above covers.
\item \textbf{Sybils.} Each principal has an independent list, so many attacker
  principals share nothing that a single one does not.
\end{itemize}

\subsection{Residual channels}
\paragraph{The public load signal.} Cross-principal hits reduce real load and so
can shift $\Lpub$. This channel carries at most $\log_2 k$ bits per epoch about
aggregate load and none about any specific prefix; $k$ and $\Delta$ are
operator knobs (\cref{sec:eval-sensitivity}).

\paragraph{Decode contention.} TBT depends on who else is decoding, as in any
shared GPU deployment, isolated caches included. \sys does not change it.

\paragraph{Mis-declared public content.} A secret registered in the public set
is unprotected. The set should hold only operator-owned or published content.

\paragraph{Outside the serving path.} Dashboards, logs or support tools that
expose physical hit rates per principal reopen the channel; \sys's guarantee
covers only the client API.
